\section{Introduction}

% [proposal] Importance of the problem
IoT deployments increasingly rely on continuous, high-rate sensor streams to drive real-time dashboards and automation. However, under bursty workloads and imperfect networks, naïvely forwarding every sensor message to the UI can overwhelm bandwidth, inflate end-to-end latency, and produce unstable or misleading displays (e.g., rapid flicker, gaps, or apparent “randomness” when updates arrive late or out of order). This is especially problematic for human-facing monitoring, where the user’s experience depends not just on whether messages eventually arrive, but on whether the displayed state remains timely, interpretable, and stable under loss, jitter, and temporary outages. Motivated by these constraints, our project studies an MQTT-based IoT gateway that can reduce communication overhead while preserving actionable dashboard behavior when network conditions degrade.

% [proposal] State of the art techniques (1/2)
A common direction in the literature is to reduce IoT communication cost by moving computation closer to the edge/fog layer: gateways and edge services filter redundant updates, aggregate multiple readings, compress payloads, or adapt sampling rates to lower network and storage pressure. Systematic reviews and mappings of edge/fog data reduction techniques emphasize that aggregation and filtering (including window-based summarization and suppression of redundant updates) are among the most frequently explored strategies because they are practical, protocol-agnostic, and provide direct reductions in transmitted volume. These approaches typically target trade-offs between data fidelity and resource usage, and they often recommend comparative evaluation under different workloads and network conditions to understand where each technique is beneficial.

% [proposal] State of the art techniques (2/2)
In parallel, MQTT has become a widely used messaging substrate for IoT systems because it offers simple publish/subscribe semantics and explicit reliability modes via QoS. In particular, QoS 0 (“at most once”) minimizes overhead but can lose messages, while QoS 1 (“at least once”) improves delivery probability at the cost of acknowledgements and possible duplicate deliveries. As a result, real systems that rely on MQTT commonly face two practical issues relevant to dashboards: (i) bursts of fine-grained updates can create unnecessary bandwidth and render jittery visualizations, and (ii) reliability mechanisms can introduce retransmissions/duplicates that must be handled at the application layer to avoid double-counting and UI instability.

% [proposal] Research gap
Despite broad interest in edge/fog data reduction, a gap remains for end-to-end, dashboard-oriented evaluation of gateway policies that jointly address (a) bandwidth pressure under bursty streams, (b) duplicates and retransmissions induced by reliability modes such as MQTT QoS 1, and (c) user-visible stability when the last hop degrades (loss/jitter/bandwidth caps/outages). Many reduction techniques are studied in isolation or under idealized conditions, and they often optimize network volume without making “what the user sees” explicit—e.g., how stale a displayed value becomes, how frequently it changes under jitter, or how tail latency behaves during impairments. Consequently, existing techniques do not fully characterize the trade-offs between compression, timeliness, and stability that determine quality for real-time monitoring.

% [proposal] How our technique bridges the gap
We propose Agrasandhani, a smart IoT gateway that treats sensor updates as a stream to be cleaned and summarized before delivery to the dashboard. The gateway combines (i) hybrid time/size batching, (ii) two-stage deduplication that drops exact duplicates by message identifier and compacts each batch to the latest update per sensor, (iii) adaptive publish-rate control that increases the batching window when publish health degrades (e.g., rising send time or queue backlog) and decreases it when conditions improve, and (iv) last-known-good caching with an explicit freshness/age indicator and TTL at the dashboard. This design directly targets the above gap by making stability and freshness first-class outcomes while retaining clear, tunable algorithmic knobs (batch window, batch cap, TTL, and adaptation thresholds) that can be evaluated systematically.

% [proposal] Experimental plan
We will implement an end-to-end MQTT pipeline (sensor simulator $\rightarrow$ gateway $\rightarrow$ dashboard subscriber/backend $\rightarrow$ web UI) and evaluate an ablation ladder of gateway variants: naive forwarding (V0), batching only (V1), batching plus compaction (V2), V2 plus adaptive publish rate (V3), and V3 plus last-known-good freshness semantics (V4). Experiments will sweep workload intensity (e.g., number of sensors and burst frequency) and MQTT reliability (QoS 0 vs QoS 1), while applying deterministic, repeatable network impairments on the gateway-to-dashboard path (bandwidth caps, loss, delay+jitter, and short outages) using standard Linux traffic control emulation. We will report end-to-end latency (including tail latency), message rate, bandwidth usage, and data loss, and we will additionally measure freshness (age of displayed values) and update jitter to capture user-visible stability under degraded networks.

% One paragraph about the results of experiments

% A list of contributions your technique/paper makes. 